% titel: Blatt 9, A8: PCl5-Dissoziation und Kc
% Teil II: der durchgerechnete Weg.

\subsection*{Gegeben und gesucht}
0.074 mol PCl$_5$(g) in 1\,L-Gefäss; bei Gleichgewicht
$[PCl_3]=0.050\ mol\,L^{-1}$. Gesucht: $[Cl_2]$, $[PCl_5]$, $K_c$.

\subsection*{Ansatz}
$PCl_5 \rightleftharpoons PCl_3 + Cl_2$. In 1 L gilt mol = mol/L.
Bei Gleichgewicht ist $x = [PCl_3] = 0.050$.

\subsection*{Rechnung}
\[
  [Cl_2] = x = 0.050\ mol\,L^{-1},
\]
\[
  [PCl_5] = 0.074 - x = 0.074-0.050 = 0.024\ mol\,L^{-1}.
\]
\[
  K_c = \frac{[PCl_3][Cl_2]}{[PCl_5]}
      = \frac{0.050\cdot0.050}{0.024}
      \approx 0.104 .
\]

\begin{ergebnis}
$[Cl_2]=0.050\ mol\,L^{-1}$, $[PCl_5]=0.024\ mol\,L^{-1}$,
$K_c \approx 0.10$.
\end{ergebnis}

\begin{kontrolle}
Dimension check: $K_c = (mol/L)^2/(mol/L) = mol/L$. Zahlen: $0.050^2/0.024
= 0.104$; passend. Da $[PCl_3]=[Cl_2]$ (vom gleichen Ursprung), stimmt die
Symmetrie der Reaktion.
\end{kontrolle}

\begin{merksatz}
Gleichgewichtskonzentration = Startkonzentration $\pm x$. Aus einer
bekannten GG-Konzentration $x$ bestimmen, Rest einsetzen, $K_c$ bilden.
\end{merksatz}
